\section*{Appendix B: Lexicon of the Collective}

\begin{description}[style=nextline, leftmargin=1em, font=\bfseries\headerfont\color{collectiveBlue}]
    \item[Aliens] The conceptual label for external market entities and legacy systems operating on capital-driven logic.
    \item[Availability Task] A scheduled task that compensates Workers for being in a "ready-to-serve" state to solve the cold-start problem for on-demand services.
    \item[Body of Knowledge (BoK)] A specialized Digital Society that maintains the global registry of Task definitions and verifies Member skill competencies.
    \item[Leaching] The strategic process of incrementally internalizing production to reduce and eventually eliminate dependency on legacy currency and external leases.
    \item[Service Node] The primary unit of production within the mesh; a recursive, self-governing entity defined by a Work Breakdown Structure (WBS).
    \item[Society] A specialized node governing the legal and territorial logic of a community instance, managing the physical "Right to Inhabit."
    \item[Unbounded Negative Balance] A feature of the Value system allowing Members to consume resources based on potential or need without punitive interest, signaled as an investment in the Member’s future labor.
\end{description}
